\subsection{Voting Theory Review}


\content{
        \begin{example} Use the Plurality method to determine the winner of this
        election. \\
        \begin{tabular}{|l|c|c|c|c|}
        \hline
        & 9 & 19 & 11 & 8 \\
        \hline
        1st & B&A&C&B\\
        \hline
        2nd & A&B&B&C\\
        \hline
        3rd & C&C&A&A\\
        \hline
        \end{tabular}
        \end{example}
}{% presentation
\vspace{2in}
}{% nopresentation
\vspace{2in}
}{% comments
}

\content{
        \begin{example} With the same preference schedule as in the previous
        example, who wins under the Pairwise Comparison method? 
        \end{example}
}{% presentation
\vspace{3in}
}{% nopresentation
\vspace{3in}\newpage
}{% comments
}

\content{
        \begin{example} Determine the winner of the election using the Borda
        Count method. \\
        \begin{tabular}{|l|c|c|c|c|}
        \hline
        & 11 & 7 & 7 & 3 \\
        \hline
        1st & A&C&B&A \\
        \hline
        2nd & C&B&C&B \\
        \hline
        3rd & B&A&A&C \\
        \hline
        \end{tabular}
        \end{example}
}{% presentation
\vspace{2in}
}{% nopresentation
\vspace{2in}
}{% comments
}

\content{
        \begin{example} Use the Plurality with Elimination to determine who wins
        this election. \\
        \begin{tabular}{|l|c|c|c|c|c|c|}
        \hline
        & 120 & 50 & 40 & 90 & 60 & 100\\
        \hline
        1st & C&B&D&A&A&D\\
        \hline
        2nd & D&C&A&C&D&B\\
        \hline
        3rd & B&A&B&B&C&A\\
        \hline
        4th & A&D&C&D&B&C\\
        \hline
        \end{tabular}
        \end{example}
}{% presentation
\vspace{3in}
}{% nopresentation
\vspace{3in}
}{% comments
}

\content{
        \begin{example} Previously, we used the Borda Count method to determine
        that the winner of this election was $C$. Did this particular election
        have a Condorcet Candidate? If so, was the Condorcet Criterion broken?
        \\
        \begin{tabular}{|l|c|c|c|c|}
        \hline
        & 11 & 7 & 7 & 3 \\
        \hline
        1st & A&C&B&A \\
        \hline
        2nd & C&B&C&B \\
        \hline
        3rd & B&A&A&C \\
        \hline
        \end{tabular}
        \end{example}
}{% presentation
\vspace{2in}
}{% nopresentation
\vspace{2in}
}{% comments
}

\content{
        \begin{example} Time to think! Does the Pairwise Comparison Method
        Satisfy the Majority Criterion? Why or why not? 
        \end{example}
}{% presentation
\vspace{2in}
}{% nopresentation
\vspace{2in}
}{% comments
}

\content{
        \begin{example} Last lecture we introduced Arrow's Impossibility
        Theorem. To demonstrate just how difficult voting can be, consider the
        following preference schedule. \\
        \begin{tabular}{|l|c|c|c|}
        \hline
        & 5 & 5 & 5 \\
        \hline
        1st & A&C&B\\
        \hline
        2nd & B&A&C \\
        \hline
        3rd & C&B&A \\
        \hline
        \end{tabular}\\
        Who should win? 
        \end{example}
}{% presentation
\vspace{2in}
}{% nopresentation
\vspace{2in}
}{% comments
}
